\documentclass[10pt]{article}
\providecommand{\FIGDIR}{fig}
\newcommand{\DIGESTDATE}{1 November 2026}
\input{preamble}

\begin{document}

% ==================================================================== MASTHEAD
\thispagestyle{empty}
\begin{tikzpicture}[remember picture, overlay]
  \fill[Deep] (current page.north west) rectangle
      ($(current page.north east)+(0,-67mm)$);
  \fill[Amber] ($(current page.north west)+(0,-67mm)$) rectangle
      ($(current page.north east)+(0,-68.4mm)$);
  \foreach \x/\y/\r in {22/12/0.9, 41/26/1.5, 63/9/0.7, 88/31/1.1, 112/17/1.9,
                        138/29/0.8, 159/13/1.3, 178/54/1.0, 196/21/1.6,
                        30/44/0.6, 74/57/1.2, 125/48/0.9, 168/6/0.8, 191/59/0.7,
                        52/18/1.0, 100/60/0.6, 148/20/1.1, 66/49/1.4, 205/45/0.8}
    \fill[Sky, opacity=0.30] ($(current page.north west)+(\x mm,-\y mm)$) circle (\r mm);
\end{tikzpicture}

\vspace*{4mm}
{\sffamily\fontsize{7.6}{9}\selectfont\color{Sky}%
 AUTOMATED DAILY LITERATURE DIGEST\;\textbullet\;PhD RESEARCH WATCH\par}
\vspace{2.2mm}
{\sffamily\bfseries\fontsize{27}{29}\selectfont\color{white}%
 Particulate Matter\par}
{\sffamily\fontsize{27}{29}\selectfont\color{Sky}%
 Monitoring \& Health Effects\par}

\vspace{5.0mm}
{\sffamily\fontsize{8.4}{10}\selectfont\color{white}%
 Issue 2026-11-01\;\;\textcolor{Amber}{\textbar}\;\;Entry window 1 November 2026 (full day)\;\;%
 \textcolor{Amber}{\textbar}\;\;7 records screened in scope\par}
\vspace{18mm}

\noindent
\begin{minipage}[t]{0.190\linewidth}\metric{7}{RECORDS IN SCOPE\\(10 EXCLUDED)}{Deep}\end{minipage}\hfill
\begin{minipage}[t]{0.190\linewidth}\metric{6}{OF 10 SUBTOPIC\\BANDS POPULATED}{Teal}\end{minipage}\hfill
\begin{minipage}[t]{0.190\linewidth}\metric{6}{EFFECT ESTIMATES\\WITH A CI}{Sky}\end{minipage}\hfill
\begin{minipage}[t]{0.190\linewidth}\metric{1}{TIER-A\\RECORDS}{Amber}\end{minipage}\hfill
\begin{minipage}[t]{0.190\linewidth}\metric{2}{RECORDS CARRIED ON\\METADATA ALONE}{Clay}\end{minipage}

\vspace{6mm}

\section*{\sffamily\bfseries\color{Deep}Signal of the day}
\vspace{-1.5mm}{\color{Amber}\rule{\linewidth}{1.1pt}}\vspace{2.5mm}

\begin{keybox}
{\sffamily\bfseries\small\color{Deep}China, 466,543 kidney deaths: coarse PM$_{2.5\text{--}10}$ per IQR (OR 1.024,
1.016--1.032) is comparable to PM$_{2.5}$ (OR 1.020, 1.012--1.027), with renal failure the most affected subtype.}\\[1.4mm]
\footnotesize
Nationwide time-stratified case-crossover, 2013--2019, modelled daily concentrations at the residential address. NO$_2$ was the
strongest of the four pollutants per IQR (+3.62\,\%, 2.47--4.78); renal-failure and glomerular deaths responded to all four,
tubulointerstitial to O$_3$ only, urolithiasis to none.
\textbf{Where it is weak:} the abstract gives no IQR widths, so the pollutant ranking is not per unit mass; lags differ by pollutant;
renal failure on a death certificate is often a proximate cause.
\textbf{Why it matters here:} a coarse-mode association this size sits in the size range where low-cost optical nodes are least reliable.\par
\end{keybox}

\vspace{3mm}

\section*{\sffamily\bfseries\color{Deep}What changed today}
\vspace{-1.5mm}{\color{Amber}\rule{\linewidth}{1.1pt}}\vspace{2.5mm}

\footnotesize
\begin{itemize}

\item \textbf{Accra schoolchildren: no pro-hypertensive signal.} 919 children, home PM$_{2.5}$ SBP $-$0.70\,mmHg ($-$3.34 to 1.94) per 10 units, in a
city at 5--8 times the WHO annual guideline. Inverse point estimates suggest confounding. LUR built on a 146-site campaign. \textsc{backfill}, entry 4 June.

\item \textbf{Low-level mortality gradient shrinks under geography.} Native American Medicare beneficiaries OR 1.12 (1.06--1.18) per $\mu$g\,m$^{-3}$ vs
1.01 (1.001--1.02); most attenuation came from state-level clustering. \textsc{backfill}, entry 4 May.

\item \textbf{SPIROMICS proteomics.} 860 plasma proteins at FDR$<$0.05 for prior-year PM$_{2.5}$, none for NO$_2$ (n$\approx$1,620, congress abstract, no effect
sizes). \textsc{backfill}, Crossref deposit 29 October.

\item \textbf{Ferrimagnetic nanoparticles in children's brains.} Mexico City autopsy tissue, 12 children and 16 adults; presence and field-driven motion shown,
no comparison group, no effect size.

\item \textbf{Sunday floor.} PubMed entered 5 candidate records; the day's Crossref yield was Elsevier deposits without abstracts. Two records ride on metadata
alone (German sick-leave claims, personal air purifier); retry the \textit{Environ Int} abstract first.

\item \textbf{Exposure basis (x1).} Of the four health records with a recoverable basis, three are modelled (national grid, LUR, prior-year surface) and one is a laboratory tissue assay: nothing measured on a person and no low-cost-sensor record today.

\end{itemize}

\vspace{2mm}

\begin{keybox}
{\sffamily\bfseries\footnotesize\color{Deep}Window continuity}\\[1.2mm]
{\fontsize{7.2}{8.8}\selectfont
The 31 October issue closed with \texttt{last\_entry\_date} at 31 October. This issue was built on the \textbf{2 November 05:0x UTC (1 November 23:0x CST)
run}; the 22:00 rule is satisfied, so 1 November is complete. It covers \textbf{1 November in full} (Sunday, first day of standard time). One record (Calder\'on-Garcidue\~nas,
EDAT 31 October) surfaced only on Europe PMC's 1 November leg and was not screened on 31 October, so it is assigned here. Three \textsc{backfill} records are carried, the working
cap, all marked with their entry dates. No rollup is owed: W45 (1--7 November) falls on Saturday 7 November.\par}
\end{keybox}

\vspace{3mm}
\par\vskip 0pt plus 18\baselineskip\penalty-200\vskip 0pt plus -18\baselineskip
\section*{\sffamily\bfseries\color{Deep}Corpus analytics}
\vspace{-1.5mm}{\color{Amber}\rule{\linewidth}{1.1pt}}\vspace{2.5mm}

\noindent\begin{minipage}{\linewidth}\includegraphics[width=\linewidth]{\FIGDIR/f1_subtopics.png}
\figcap{Subtopic assignment of the 7 records across six bands: burden and policy two (sick leave, Medicare mortality); respiratory,
cardiovascular, neuro, occupational \& indoor and other clinical one each. Mechanistic, reproductive, exposure and sensing bands are empty.}\end{minipage}

\vspace{2mm}
\noindent\begin{minipage}{\linewidth}\includegraphics[width=\linewidth]{\FIGDIR/f2_design_endpoint.png}
\figcap{Left --- architecture: observational cross-sectional three (Accra, Medicare case-control, SPIROMICS baseline), metadata-only two, and one each
acute case-crossover and laboratory (autopsy tissue). Right --- endpoints from kidney mortality and blood pressure to the plasma proteome and sick leave.}\end{minipage}

\vspace{2mm}
\noindent\begin{minipage}{\linewidth}\includegraphics[width=\linewidth]{\FIGDIR/f3_metric_geography.png}
\figcap{Left --- particle metric: PM$_{2.5}$-only 4, unspeciated 3. Right --- geography: North America 2, then one each Sub-Saharan Africa, Europe, Latin
America and China; the personal-purifier record states no site and is drawn under global/multi-region.}\end{minipage}

\vspace{2mm}
\noindent\begin{minipage}{\linewidth}\includegraphics[width=\linewidth]{\FIGDIR/f6_lifecourse.png}
\figcap{Life-course windows: two childhood records (Accra blood pressure, Mexico City brain tissue), one working-age (sick leave, metadata only) and one
older-adult (Medicare mortality). No in-utero or separate adolescent record.}\end{minipage}

\vspace{2mm}
\noindent\begin{minipage}{\linewidth}\includegraphics[width=\linewidth]{\FIGDIR/f4_heatmap.png}
\figcap{Subtopic $\times$ study architecture. Seven records occupy seven cells; no cell has a replicate today, so this is a map of what arrived, not a
pattern.}\end{minipage}

\vspace{2mm}

\IfFileExists{\FIGDIR/extra_figs.tex}{\input{\FIGDIR/extra_figs.tex}}{}



\noindent\begin{minipage}{\linewidth}\includegraphics[width=\linewidth]{\FIGDIR/f5_forest.png}
\figcap{Ratio scale, five estimates from three studies. Increments differ (per IQR with undisclosed widths in China; per 1\,$\mu$g\,m$^{-3}$ in Medicare; per 10
units in Accra) and the Medicare slope is a low-decile contrast: nothing here is comparable across studies or poolable.}\end{minipage}

\vspace{2mm}
\noindent\begin{minipage}{\linewidth}\includegraphics[width=\linewidth]{\FIGDIR/f5_forest_2.png}
\figcap{Difference scale, one estimate: systolic blood pressure per 10 units of home PM$_{2.5}$, Accra schoolchildren (n=919), an inverse point estimate with a CI spanning zero.}\end{minipage}

\vspace{4mm}

\par\vskip 0pt plus 24\baselineskip\penalty-200\vskip 0pt plus -24\baselineskip
\section*{\sffamily\bfseries\color{Deep}Paper-level digest}
\vspace{-1.5mm}{\color{Amber}\rule{\linewidth}{1.1pt}}\vspace{2.5mm}

{\fontsize{7.4}{9}\selectfont\color{Slate}Ordered by cluster. Each entry gives the
methodological core and the finding that matters, not an abstract paraphrase. Records
marked \textit{metadata only} carry no recoverable abstract and assert no finding.\par}

\band{Neuro / mental health\hfill 1 record}{Deep}

\paperentry{Calderon-Garciduenas ferrimagnetic nanoparticles olfactory bulb cortex Mexico City children 2026}
{The olfactory bulb, prefrontal and anterior cingulate cortex are early pediatric targets for anthropogenic magnetic nanoparticles and their motion behavior under electromagnetic fields.}
{Gac Med Mex}
{48 brain samples (olfactory bulb, anterior cingulate, prefrontal cortex) from 12 children (8.5 +/- 7.6 y) and 16 adults (39.2 +/- 19.6 y), Metropolitan Mexico City. Magnetophoresis at 25-1,000 mT AC showed small displacements and significant vector-direction variation; TEM and EDX located ferrimagnetic nanoparticles. No effect size. Limitation: autopsy convenience series with no comparison population or exposure measure in the abstract, so 'early pediatric targets' and the neural-damage risk are inferred, not tested; particle attribution to anthropogenic sources is not apportioned here. Relevance: low for sensing - mass-calibrated optical nodes are blind to this metallic ultrafine fraction, which is the actual argument for co-located number or magnetic measurements.}
{10.24875/gmm.m26001158}

\par\vskip 0pt plus 10\baselineskip\penalty-200\vskip 0pt plus -10\baselineskip
\band{Cardiovascular \& metabolic\hfill 1 record}{Teal}

\paperentry{Alli ambient air pollution blood pressure schoolchildren Accra LUR 2026}
{Urban ambient air pollution and blood pressure in schoolchildren in Accra, Ghana}
{Sci Rep}
{BACKFILL (entry date 4 Jun 2026; Amass). 919 children aged 7-14 in 90 schools; one-third with elevated BP. Per 10 units home exposure: SBP PM$_{2.5}$ -0.70 mmHg (-3.34 to 1.94), BC -4.61 (-9.59 to 0.37); DBP PM$_{2.5}$ -0.83 (-2.93 to 1.28), BC -4.35 (-8.19 to -0.50); elevated BP PM$_{2.5}$ OR 0.75 (0.39-1.41). Exposures ran 5-8 times the WHO annual PM$_{2.5}$ guideline. Limitation: cross-sectional, and inverse associations in a high-exposure city suggest residual confounding by socioeconomic or adiposity patterning; the upper bound does not exclude a moderate SBP rise. Relevance: high - the LUR rests on a 146-site campaign in sub-Saharan Africa, the design a low-cost network could sustain.}
{10.1038/s41598-026-54135-6}

\par\vskip 0pt plus 10\baselineskip\penalty-200\vskip 0pt plus -10\baselineskip
\band{Respiratory \& allergic\hfill 1 record}{Sky}

\paperentry{Ruvuna ambient PM$_{2.5}$ NO$_2$ plasma proteomics SPIROMICS SomaScan congress 2026}
{Associations of ambient air pollution and protein expression among individuals from the Subpopulations and intermediate outcome measures in COPD study (SPIROMICS)}
{Eur Respir J (congress abstract)}
{BACKFILL (Crossref deposit 29 Oct 2026; Amass; no PMID). ERS congress abstract. SPIROMICS tobacco-exposed cohort, n=1,627 for PM$_{2.5}$ and n=1,618 for NO$_2$; SomaScan 7k against modelled prior-year exposure, adjusted for site, COPD status, smoking, BMI, occupation, income and deprivation. 860 proteins associated with PM$_{2.5}$ at FDR<0.05, none with NO$_2$; enrichment in vascular remodelling, insulin signalling and inflammation. Limitation: abstract only, no effect sizes or replication, cross-sectional, and a one-pollutant-sided result may partly reflect exposure-model quality or spatial confounding. Relevance: low - no sensing element, but the proteomic contrast depends on exposure accuracy.}
{10.1183/13993003.congress-2026.oa5535}

\par\vskip 0pt plus 10\baselineskip\penalty-200\vskip 0pt plus -10\baselineskip
\band{Occupational \& indoor\hfill 1 record}{Amber}

\paperentry{Nie personal air purification device multi-scale modelling fan speed occupant exposure 2026}
{Multi-scale modeling and fan speed optimization of a personal air-purification device for occupant exposure control in polluted indoor environments}
{Build Environ}
{\textsc{metadata only.} Metadata only: Elsevier day-0 deposit, no abstract in Crossref or Semantic Scholar; no finding asserted. The title describes multi-scale modelling and fan-speed optimisation of a personal air-purification device. Limitation: the title does not confirm particulate matter as the pollutant, and no performance figures, validation or occupant data are available. Relevance: to be assessed - breathing-zone control is where stationary sensors misstate exposure; retain only if the abstract confirms a PM target and measured validation.}
{10.1016/j.buildenv.2026.115523}

\par\vskip 0pt plus 10\baselineskip\penalty-200\vskip 0pt plus -10\baselineskip
\band{Burden, policy \& mitigation\hfill 2 records}{Coral}

\paperentry{Aretz sick leave long-term fine particulate German health claims 13-year 2026}
{Assessing the health and economic burden of sick leave associated with long-term fine particulate matter exposure: a 13-year analysis using German health claims data}
{Environ Int}
{\textsc{metadata only.} Metadata only: Elsevier day-0 deposit, no abstract in Crossref, Europe PMC or Semantic Scholar; no finding asserted. The title describes a 13-year analysis of German health claims linking long-term fine PM to sick leave and its economic burden. Limitation: exposure source, increment, cohort size and estimate are all unknown until the abstract deposits. Relevance: to be assessed - claims-linked sick leave is a working-age endpoint rarely paired with PM, and this is the first record to retry on the next weekly abstract sweep.}
{10.1016/j.envint.2026.110586}

\paperentry{Hess low-level PM$_{2.5}$ mortality Native American Medicare geographic attenuation 2026}
{Untangling Impacts of Socioeconomic Position, Chronic Disease, and Low-Level PM$_{2.5}$ Exposure on Mortality Among Native American Medicare Beneficiaries}
{Int J Environ Res Public Health}
{BACKFILL (entry date 4 May 2026; Amass). Case-control within 1,713,399 traditional Medicare beneficiaries in the lowest PM$_{2.5}$ decile, 2015-2016, with PCA summaries of area SEP and chronic disease. Fully adjusted OR per $\mu$g\,m$^{-3}$: Native American 1.12 (1.06-1.18), non-Native American 1.01 (1.001-1.02). Sequential adjustment attributed most attenuation in Native American beneficiaries to state-level clustering, less to SEP and disease. Limitation: the abstract gives no pre-adjustment coefficient, zip-level exposure over a narrow range makes a per-$\mu$g\,m$^{-3}$ slope fragile, and state effects may absorb real exposure contrast. Relevance: moderate - low-level gradients are where sparse monitoring error is largest.}
{10.3390/ijerph23040464}

\par\vskip 0pt plus 10\baselineskip\penalty-200\vskip 0pt plus -10\baselineskip
\band{Other clinical endpoints\hfill 1 record}{Slate}

\paperentry{Xiao air pollution kidney disease mortality subtypes China case-crossover 2026}
{Air pollution exposure and death from different subtypes of kidney diseases: A nationwide, individual-level, case-crossover study in China.}
{Chin Med J (Engl)}
{466,543 kidney deaths, China 2013-2019, time-stratified case-crossover with conditional logistic regression. Per IQR: PM$_{2.5}$ lag 0-2 OR 1.0195 (1.0122-1.0269); coarse PM$_{2.5}$-10 lag 0-3 OR 1.0240 (1.0158-1.0323); O$_3$ +2.05\% (1.05-3.07); NO$_2$ lag 0-3 +3.62\% (2.47-4.78). Renal failure and glomerular deaths responded to all four pollutants; tubulointerstitial deaths to O$_3$ only; urolithiasis null. Limitation: the abstract does not give the IQR widths, so 'NO$_2$ strongest' is not a per-mass comparison, the lag differs by pollutant, and renal failure on a death certificate is often a proximate cause. Relevance: moderate - coarse-mode signal comparable to PM$_{2.5}$ per IQR, in a size fraction low-cost optical sensors undercount.}
{10.1097/cm9.0000000000004318}

\clearpage
\section*{\sffamily\bfseries\color{Deep}Corpus register}
\vspace{-1.5mm}{\color{Amber}\rule{\linewidth}{1.1pt}}\vspace{2.5mm}

\input{table_rows}

\vspace{2mm}

\begin{keybox}
{\sffamily\bfseries\footnotesize\color{Deep}Provenance and method}\\[1.2mm]
{\fontsize{7.2}{8.8}\selectfont
Entry window \textbf{1 November 2026}, single day (Sunday). Harvester legs: \textbf{PubMed} health \textbf{0}, sensing \textbf{1}; \textbf{Europe PMC}
\textbf{9} after dropping 2 AGRICOLA retro-index records; \textbf{Crossref by ISSN} \textbf{6}; \textbf{arXiv} 150 kB, no in-window entry; \textbf{OpenAlex} HTTP 429 (the OpenAlex connector
is not on this session). 14 unique, 6 already carried, 8 fresh; PubMed connector (EDAT 1 Nov, three axes) 5 PMIDs, 1 seen, 4 appended (\textbf{12 screened}). Abstracts: 4 recovered by efetch; Elsevier
day-0 deposits remain empty. \textbf{Amass} BiomedCore (ingest date filter) surfaced \textbf{3 backfill} records admitted and 2 rejected; others held. \textbf{Consensus} calibration sweep: 10 hits, 0 new (Fang 2026 already shipped; Feng 2026 review already rejected).
\textbf{medRxiv} toolbox leg: no keyword filter, no PM record. bioRxiv/medRxiv and alphaXiv connectors are not on this session. \textbf{Scite} retraction check unavailable (connector requires a paid plan); Amass reports no retraction on PMID-bearing backfills. \textbf{Wiley Scholar Gateway} not called (needs re-authorisation).
\textbf{7 admitted} (4 window, 3 backfill), \textbf{10 rejected} with logged reasons. \textbf{Trial watch} 1--2 November: ClinicalTrials.gov 0; Amass TrialCore 10 returned, one relevant (NCT07311967, UIC cooking UFP chamber, not previously logged). \textbf{2 carried on metadata alone}.
Abstracts remain the copyright of their publishers.\par}
\end{keybox}

\end{document}
